\chapter{Chapter02 — Degree One Rigidity}
%%% generated by the blueprint pipeline from FABL.Chapter02.InfluencesAndDerivatives.DegreeOneRigidity
\section{Overview}

Declarations of the Lean module \texttt{FABL.Chapter02.InfluencesAndDerivatives.DegreeOneRigidity}.

\section{Declarations}

\begin{theorem}[Fourier coeff eq zero of fourier weight at level one eq one]
\lean{FABL.fourierCoeff_eq_zero_of_fourierWeightAtLevel_one_eq_one}
\statementsource{Analysis-of-Boolean-Functions-by-Ryan-ODonnell}{
  pdf-ffbe9c7b3306-p025-b006,
  pdf-ffbe9c7b3306-p037-b008
}
\label{FABL.fourierCoeff_eq_zero_of_fourierWeightAtLevel_one_eq_one}
\leanok
\uses{FABL.BooleanFunction, FABL.BooleanFunction.toReal, FABL.fourierCoeff, FABL.fourierWeight, FABL.fourierWeightAtLevel, FABL.sum_sq_fourierCoeff_eq_one}
\difficulty{4}
% dataset-pipeline: stub -- prose pending build_informal_statements.py
Statement pending regeneration.
\end{theorem}

\begin{theorem}[Level one Fourier expansion]
\lean{FABL.toReal_eq_sum_singleton_of_fourierWeightAtLevel_one_eq_one}
\label{FABL.toReal_eq_sum_singleton_of_fourierWeightAtLevel_one_eq_one}
\leanok
\uses{FABL.BooleanFunction, FABL.BooleanFunction.toReal, FABL.fourierCoeff, FABL.fourierCoeff_eq_zero_of_fourierWeightAtLevel_one_eq_one, FABL.fourierWeightAtLevel, FABL.fourier_expansion, FABL.monomial, FABL.signValue}
\difficulty{4}
Let $f$ be a function from the $n$-dimensional sign cube to $\{-1,1\}$, and suppose the
total Fourier weight of its level-one coefficients is $1$. Then at every input $x$, its
real-valued form satisfies $f(x)=\sum_{i=1}^n \widehat f(\{i\})x_i$, where $\widehat
f(\{i\})$ is the Fourier coefficient indexed by the singleton coordinate $i$.
\end{theorem}

\begin{theorem}[Derivative equals singleton Fourier coefficient]
\lean{FABL.discreteDerivative_toReal_eq_singletonCoeff_of_fourierWeightAtLevel_one_eq_one}
\label{FABL.discreteDerivative_toReal_eq_singletonCoeff_of_fourierWeightAtLevel_one_eq_one}
\leanok
\uses{FABL.BooleanFunction, FABL.BooleanFunction.toReal, FABL.discreteDerivative, FABL.discreteDerivative_eq_fourier_sum, FABL.fourierCoeff, FABL.fourierCoeff_eq_zero_of_fourierWeightAtLevel_one_eq_one, FABL.fourierWeightAtLevel, FABL.monomial}
\difficulty{4}
Let $f$ be a function from the $n$-dimensional sign cube to $\{\pm1\}$, and suppose the
sum of the squared Fourier coefficients of its real-valued version over all singleton
subsets is $1$. Then, for every coordinate $i$ and every cube point $x$, the discrete
derivative of the real-valued version of $f$ in coordinate $i$ at $x$ equals its Fourier
coefficient on the singleton $\{i\}$.
\end{theorem}

\begin{theorem}[Level-One Fourier Weight]
\lean{FABL.fourierWeightAtLevel_one_eq_sum_singleton}
\label{FABL.fourierWeightAtLevel_one_eq_sum_singleton}
\leanok
\uses{FABL.fourierCoeff, FABL.fourierWeight, FABL.fourierWeightAtLevel}
\difficulty{4}
For a function $f:\{-1,1\}^n\to\mathbb{R}$ on the $n$-dimensional Boolean cube, its
Fourier weight at level one is the sum of the squares of its singleton Fourier
coefficients: $W_1(f)=\sum_{i=1}^n \widehat f(\{i\})^2$.
\end{theorem}

\begin{theorem}[Expansion of the Degree-One Fourier Weight]
\lean{FABL.fourierWeightAtMost_one_eq_empty_add_sum_singleton}
\label{FABL.fourierWeightAtMost_one_eq_empty_add_sum_singleton}
\leanok
\uses{FABL.fourierCoeff, FABL.fourierWeight, FABL.fourierWeightAtMost}
\difficulty{5}
For any real-valued function $f$ on the $n$-dimensional sign cube, the sum of its
squared Fourier coefficients over subsets of size at most one equals the square of its
coefficient at the empty subset plus the sum, over all coordinates, of the squared
coefficients at the corresponding singleton subsets.
\end{theorem}

\begin{theorem}[Singleton coeff sq eq zero or one of fourier weight at level one eq one]
\lean{FABL.singletonCoeff_sq_eq_zero_or_one_of_fourierWeightAtLevel_one_eq_one}
\statementsource{Analysis-of-Boolean-Functions-by-Ryan-ODonnell}{
  pdf-ffbe9c7b3306-p025-b006,
  pdf-ffbe9c7b3306-p028-b002,
  pdf-ffbe9c7b3306-p037-b008
}
\label{FABL.singletonCoeff_sq_eq_zero_or_one_of_fourierWeightAtLevel_one_eq_one}
\leanok
\uses{FABL.BooleanFunction, FABL.BooleanFunction.toReal, FABL.IsPivotal, FABL.discreteDerivative_toReal_eq_singletonCoeff_of_fourierWeightAtLevel_one_eq_one, FABL.fourierCoeff, FABL.fourierWeightAtLevel, FABL.pivotalIndicator, FABL.sq_discreteDerivative_toReal_eq_pivotalIndicator}
\difficulty{4}
% dataset-pipeline: stub -- prose pending build_informal_statements.py
Statement pending regeneration.
\end{theorem}

\begin{theorem}[Concentration of level-one Fourier weight]
\lean{FABL.exists_unique_singletonCoeff_sq_eq_one}
\label{FABL.exists_unique_singletonCoeff_sq_eq_one}
\leanok
\uses{FABL.BooleanFunction, FABL.BooleanFunction.toReal, FABL.fourierCoeff, FABL.fourierWeightAtLevel, FABL.fourierWeightAtLevel_one_eq_sum_singleton, FABL.singletonCoeff_sq_eq_zero_or_one_of_fourierWeightAtLevel_one_eq_one}
\difficulty{5}
Let $f$ be a function from the $n$-dimensional sign cube to $\{ -1,1 \}$. If the sum
of the squares of its level-one Fourier coefficients is $1$, then there is a coordinate
$i$ whose Fourier coefficient has square $1$, and every other level-one Fourier
coefficient is zero.
\end{theorem}

\begin{theorem}[Eq dictator or neg dictator of fourier weight at level one eq one]
\lean{FABL.eq_dictator_or_neg_dictator_of_fourierWeightAtLevel_one_eq_one}
\statementsource{Analysis-of-Boolean-Functions-by-Ryan-ODonnell}{
  pdf-ffbe9c7b3306-p037-b008,
  pdf-ffbe9c7b3306-p059-b004
}
\label{FABL.eq_dictator_or_neg_dictator_of_fourierWeightAtLevel_one_eq_one}
\leanok
\uses{FABL.BooleanFunction, FABL.BooleanFunction.toReal, FABL.dictator, FABL.exists_unique_singletonCoeff_sq_eq_one, FABL.fourierCoeff, FABL.fourierWeightAtLevel, FABL.signValue, FABL.signValue_injective, FABL.toReal_eq_sum_singleton_of_fourierWeightAtLevel_one_eq_one}
\difficulty{4}
% dataset-pipeline: stub -- prose pending build_informal_statements.py
Statement pending regeneration.
\end{theorem}

\begin{theorem}[Fourier coeff eq zero of fourier weight at most one eq one]
\lean{FABL.fourierCoeff_eq_zero_of_fourierWeightAtMost_one_eq_one}
\statementsource{Analysis-of-Boolean-Functions-by-Ryan-ODonnell}{
  pdf-ffbe9c7b3306-p025-b006,
  pdf-ffbe9c7b3306-p028-b002,
  pdf-ffbe9c7b3306-p037-b008
}
\label{FABL.fourierCoeff_eq_zero_of_fourierWeightAtMost_one_eq_one}
\leanok
\uses{FABL.BooleanFunction, FABL.BooleanFunction.toReal, FABL.fourierCoeff, FABL.fourierWeight, FABL.fourierWeightAtMost, FABL.sum_sq_fourierCoeff_eq_one}
\difficulty{4}
% dataset-pipeline: stub -- prose pending build_informal_statements.py
Statement pending regeneration.
\end{theorem}

\begin{theorem}[Derivative equals singleton Fourier coefficient]
\lean{FABL.discreteDerivative_toReal_eq_singletonCoeff_of_fourierWeightAtMost_one_eq_one}
\label{FABL.discreteDerivative_toReal_eq_singletonCoeff_of_fourierWeightAtMost_one_eq_one}
\leanok
\uses{FABL.BooleanFunction, FABL.BooleanFunction.toReal, FABL.discreteDerivative, FABL.discreteDerivative_eq_fourier_sum, FABL.fourierCoeff, FABL.fourierCoeff_eq_zero_of_fourierWeightAtMost_one_eq_one, FABL.fourierWeightAtMost, FABL.monomial}
\difficulty{4}
Let $f$ be a function from the $n$-dimensional sign cube to $\{\pm1\}$, and suppose the
sum of the squared Fourier coefficients of its real-valued version over all subsets of
size at most one is $1$. Then, for every coordinate $i$ and every cube point $x$, the
discrete derivative of the real-valued version of $f$ in coordinate $i$ at $x$ equals
its Fourier coefficient on the singleton $\{i\}$.
\end{theorem}

\begin{theorem}[Singleton coeff sq eq zero or one of fourier weight at most one eq one]
\lean{FABL.singletonCoeff_sq_eq_zero_or_one_of_fourierWeightAtMost_one_eq_one}
\statementsource{Analysis-of-Boolean-Functions-by-Ryan-ODonnell}{
  pdf-ffbe9c7b3306-p025-b006,
  pdf-ffbe9c7b3306-p028-b002,
  pdf-ffbe9c7b3306-p037-b008
}
\label{FABL.singletonCoeff_sq_eq_zero_or_one_of_fourierWeightAtMost_one_eq_one}
\leanok
\uses{FABL.BooleanFunction, FABL.BooleanFunction.toReal, FABL.IsPivotal, FABL.discreteDerivative_toReal_eq_singletonCoeff_of_fourierWeightAtMost_one_eq_one, FABL.fourierCoeff, FABL.fourierWeightAtMost, FABL.pivotalIndicator, FABL.sq_discreteDerivative_toReal_eq_pivotalIndicator}
\difficulty{4}
% dataset-pipeline: stub -- prose pending build_informal_statements.py
Statement pending regeneration.
\end{theorem}

\begin{theorem}[Is k junta one of fourier weight at most one eq one]
\lean{FABL.isKJunta_one_of_fourierWeightAtMost_one_eq_one}
\statementsource{Analysis-of-Boolean-Functions-by-Ryan-ODonnell}{
  pdf-ffbe9c7b3306-p037-b008,
  pdf-ffbe9c7b3306-p044-b003
}
\label{FABL.isKJunta_one_of_fourierWeightAtMost_one_eq_one}
\leanok
\uses{FABL.BooleanFunction, FABL.BooleanFunction.toReal, FABL.IsKJunta, FABL.dictator, FABL.eq_dictator_or_neg_dictator_of_fourierWeightAtLevel_one_eq_one, FABL.fourierCoeff, FABL.fourierCoeff_eq_zero_of_fourierWeightAtMost_one_eq_one, FABL.fourierWeightAtLevel, FABL.fourierWeightAtLevel_one_eq_sum_singleton, FABL.fourierWeightAtMost, FABL.fourier_expansion, FABL.monomial, FABL.signValue_injective, FABL.singletonCoeff_sq_eq_zero_or_one_of_fourierWeightAtMost_one_eq_one, FABL.sum_sq_fourierCoeff_eq_one}
\difficulty{5}
% dataset-pipeline: stub -- prose pending build_informal_statements.py
Statement pending regeneration.
\end{theorem}
